% Wave-17 Opus-13: full first-principles derivation of the 6D holomorphic
% Chern--Simons avatar as an E_3-factorisation algebra in the
% Costello--Francis--Gwilliam 2026 machine. Author: Raeez Lorgat.
% Chriss--Ginzburg voice. Adversarial-heal output; derivable from first
% principles side-by-side with CFG26.

\providecommand{\ClaimStatusTheorem}{\textsc{[theorem]}}
\providecommand{\ClaimStatusProposition}{\textsc{[proposition]}}
\providecommand{\ClaimStatusLemma}{\textsc{[lemma]}}
\providecommand{\ClaimStatusCorollary}{\textsc{[corollary]}}
\providecommand{\ClaimStatusDefinition}{\textsc{[definition]}}
\providecommand{\ClaimStatusDerivation}{\textsc{[first-principles]}}
\providecommand{\ClaimStatusRemark}{\textsc{[remark]}}
\providecommand{\ClaimStatusConjectured}{\textsc{[conjectured]}}

\providecommand{\CC}{\mathbb{C}}
\providecommand{\RR}{\mathbb{R}}
\providecommand{\ZZ}{\mathbb{Z}}
\providecommand{\QQ}{\mathbb{Q}}
\providecommand{\fg}{\mathfrak{g}}
\providecommand{\cA}{\mathcal{A}}
\providecommand{\cB}{\mathcal{B}}
\providecommand{\cE}{\mathcal{E}}
\providecommand{\cF}{\mathcal{F}}
\providecommand{\cO}{\mathcal{O}}
\providecommand{\cH}{\mathcal{H}}
\providecommand{\cK}{\mathcal{K}}
\providecommand{\cM}{\mathcal{M}}
\providecommand{\cD}{\mathcal{D}}
\providecommand{\cC}{\mathcal{C}}
\providecommand{\hCS}{\mathrm{hCS}}
\providecommand{\Conf}{\mathrm{Conf}}
\providecommand{\Obs}{\mathrm{Obs}}
\providecommand{\CE}{\mathrm{CE}}
\providecommand{\HH}{\mathrm{HH}}
\providecommand{\Sym}{\mathrm{Sym}}
\providecommand{\Kos}{\mathsf{K}}
\providecommand{\chir}{\mathrm{chir}}
\providecommand{\Dolb}{\mathrm{Dolb}}
\providecommand{\Tot}{\mathrm{Tot}}
\providecommand{\hbarb}{\hbar}
\providecommand{\BV}{\mathrm{BV}}
\providecommand{\BRST}{\mathrm{BRST}}

\section{Six-dimensional holomorphic Chern--Simons as a Costello--Francis--Gwilliam factorisation algebra: first-principles derivation}
\label{opus13:sec:main}

\subsection*{Prolegomenon}

The programme's $\mathbb E_3$-factorisation side is anchored by a single
quantum field theory: six-dimensional holomorphic Chern--Simons on a
Calabi--Yau threefold. Its classical BV datum is three lines; its
quantum observable complex is the Costello--Gwilliam factorisation
algebra with Bochner--Martinelli propagator; its one-loop obstruction
is the quartic adjoint Casimir integrated against $c_3(X)$; its
deformation modulus is the Chevalley--Eilenberg realisation of the
$\mathbb E_3$-Hochschild complex in degree~$3$; its dualisability
reduces to the $\mathbb E_3$-Koszul self-duality $\cD_3^! \simeq
\mathrm{Lie}[2]$ realised in the pro-completed homotopy operad.

Every line below is derivable from primary first principles, with
Feynman diagrams, Bochner--Martinelli integrands, and Koszul signs
written out explicitly. The chapter reads against
Costello--Francis--Gwilliam~2026 side by side; at every step the reader
can produce the matching CFG derivation by hand.

Conventions: $X$ a Calabi--Yau threefold (so $\Omega_X \in H^{3,0}(X)$
nonvanishing), $\fg$ a finite-dimensional Lie algebra over $\CC$ with
invariant pairing $\kappa_\fg$, and $\hbar$ the deformation parameter
of the BV quantisation. Cohomological degree (ghost degree + form
degree) is denoted $|\cdot|$; $[a,b]$ denotes the graded Lie bracket
with Koszul sign $(-1)^{|a||b|}$ on swap. Dolbeault operator
$\bar\partial$ has degree $+1$; de Rham exterior derivative on form
degree composes with cohomological shift.


\subsection{Classical BV datum}
\label{opus13:subsec:classical}

\subsubsection{Field space, ghost structure, action}

\begin{definition}[6d hCS BV field]\ClaimStatusDefinition
\label{opus13:def:bv-field}
Fix a CY$_3$ $X$ with holomorphic volume $\Omega_X$. The BV field is
\[
  \cA \;=\; c \;+\; A^{(0,1)} \;+\; A^{*(0,2)} \;+\; c^{*(0,3)} \;\in\;
  \Omega^{0,\bullet}(X, \fg)[1],
\]
where the cohomological shift by $[1]$ means the component of Dolbeault
degree $q$ sits in cohomological degree $q - 1$: $c$ (ghost) in degree
$-1$, $A^{(0,1)}$ (gauge field) in degree $0$, $A^{*(0,2)}$ (anti-field)
in degree $+1$, $c^{*(0,3)}$ (anti-ghost) in degree $+2$.
\end{definition}

\begin{proposition}[Classical BV action]\ClaimStatusProposition
\label{opus13:prop:cl-action}
The classical BV action is
\begin{equation}
  S_{\mathrm{cl}}(\cA) \;=\; \int_X \Omega_X \wedge
  \bigl\langle \cA,\; \bar\partial \cA + \tfrac{1}{3}[\cA, \cA]\bigr\rangle
  \label{opus13:eq:S-cl}
\end{equation}
with $\langle-,-\rangle$ the $\fg$-invariant pairing combined with the
natural pairing $\Omega^{0,p}(X,\fg) \otimes \Omega^{0,q}(X,\fg)\to
\Omega^{0,p+q}(X,\CC)$. The integrand is nonzero only on the
$(0,3)$-component of $\Omega_X \wedge (\cdots)$, where it is a top form.
\end{proposition}

\begin{proof}
Expand $\cA = \sum_{q=0}^{3} \cA^{(0,q)}$ with $\cA^{(0,q)}$ the
cohomological-degree-$(q-1)$ component, and insert into
\eqref{opus13:eq:S-cl}. The quadratic kinetic term
$\int_X \Omega_X \wedge \langle \cA,\bar\partial\cA\rangle$ selects, via
the bi-grading filter on $\Omega_X\wedge\Omega^{0,\bullet}\wedge
\Omega^{0,\bullet}$, the pairings $\langle A^{(0,1)},\bar\partial
A^{(0,1)}\rangle$ (yielding the abelian kinetic bilinear on the gauge
field), $\langle c,\bar\partial A^{*(0,2)}\rangle$ (the
ghost--anti-field BRST pairing), $\langle A^{*(0,2)},\bar\partial c
\rangle$ (its transpose), and $\langle c,\bar\partial c^{*(0,3)}\rangle$
(ghost--anti-ghost). Integration by parts and the Jacobi identity
collapse pairs; the cubic term
$\int_X \Omega_X \wedge \tfrac{1}{3}\langle\cA,[\cA,\cA]\rangle$
expands by the antisymmetric Koszul rule
$[a,b] = -(-1)^{|a||b|}[b,a]$ and reduces to the graded-cyclic
$\langle a,[b,c]\rangle = \langle b,[c,a]\rangle(-1)^{|a|(|b|+|c|)}$.
\end{proof}

\begin{remark}[Gauge-invariance and master equation]\ClaimStatusRemark
\label{opus13:rem:mastereq}
The classical master equation $\{S_{\mathrm{cl}}, S_{\mathrm{cl}}\} = 0$
follows from the Jacobi identity on $\fg$ together with
$\bar\partial^2 = 0$ on $X$ and the CY property
$d\Omega_X = \bar\partial\Omega_X = 0$. The antibracket
$\{-,-\}$ is the $(-1)$-shifted Poisson bracket induced by the BV
symplectic pairing~\eqref{opus13:eq:sympl} below.
\end{remark}

\subsubsection{\texorpdfstring{$(-1)$-shifted symplectic structure via Serre duality}{(-1)-shifted symplectic structure via Serre duality}}

\begin{proposition}[BV symplectic pairing]\ClaimStatusProposition
\label{opus13:prop:sympl}
The BV pairing
\begin{equation}
  \omega_{\BV}(\cA_1, \cA_2) \;=\; \int_X \Omega_X \wedge
  \langle \cA_1, \cA_2\rangle
  \label{opus13:eq:sympl}
\end{equation}
is a $(-1)$-shifted symplectic form on the space of fields
$\Omega^{0,\bullet}(X,\fg)[1]$.
\end{proposition}

\begin{proof}
Serre duality on $X$:
$\Omega^{0,q}(X,\fg)^\vee \simeq \Omega^{0,3-q}(X,\fg^\vee)$ via pairing
with $\Omega_X \wedge (-)$ and the invariant pairing on $\fg$. A form
$\cA^{(0,q)}$ of cohomological degree $q - 1$ pairs with
$\cA^{(0,3-q)}$ of cohomological degree $2 - q$ to give total
cohomological degree $-1$. The antisymmetry
$\omega_{\BV}(\cA_1,\cA_2) = -(-1)^{|\cA_1||\cA_2|}\omega_{\BV}(\cA_2,\cA_1)$
comes from graded cyclicity of $\langle-,-\rangle$ plus the wedge
identity $\alpha\wedge\beta = (-1)^{|\alpha||\beta|}\beta\wedge\alpha$
for Dolbeault forms. Closure under the BV differential follows from
CY-triviality $\bar\partial\Omega_X = 0$ plus Stokes.
\end{proof}

\subsubsection{AKSZ shift law and the dimensional ladder}

\begin{theorem}[AKSZ shift law, PTVV]\ClaimStatusTheorem
\label{opus13:thm:aksz-shift}
Let $Y$ be a Calabi--Yau $d$-fold and consider AKSZ--Chern--Simons with
target $B\fg$. The moduli space $\mathrm{Map}(Y_{\DR}, B\fg)_{\fg\neq 0}$
carries an $(3-2d)$-shifted symplectic structure. After BV completion
(pre-imposing ghosts and anti-fields) the BV complex of fields on
$\Omega^{0,\bullet}(Y,\fg)[1]$ has a $(-1)$-shifted symplectic pairing
for every $d$, consistent via AKSZ descent:
\begin{center}
\begin{tabular}{@{}llll@{}}
$d$ & AKSZ moduli shift & AKSZ $\mathbb E_n$-tangent & BV pairing \\
\hline
$1$ & $+1$-shifted & $\mathbb E_3$-Poisson & $(-1)$-shifted symplectic \\
$2$ & $-1$-shifted & $\mathbb E_1$ (Lie--algebra-like) & $(-1)$-shifted symplectic \\
$3$ & $-3$-shifted & $\mathbb E_1$-\textit{chiral} on a curve, $\mathbb E_3$-hFA on $\CC^3$ & $(-1)$-shifted symplectic \\
$4$ & $-5$-shifted & $\mathbb E_5$-Poisson (not symplectic) & $(-1)$-shifted symplectic
\end{tabular}
\end{center}
In every case the BV pairing on fields is $(-1)$-shifted; what varies
across the ladder is the AKSZ moduli shift and hence the Lurie
centre / $\mathbb E_n$-dimensionality inherited at the dg level.
Source: Pantev--Toen--Vaquie--Vezzosi 2013 \emph{Publ. IHES} 117
Thm~2.5; Costello--Gwilliam Vol I §3.6.
\end{theorem}

\begin{remark}[Why the source prompt's ``shift = d-4'' does not parse literally]\ClaimStatusRemark
\label{opus13:rem:shift-law-disambiguation}
The formula $\mathrm{shift} = d - 4$ in the source prompt conflates two
distinct shifts. (i) The \emph{BV symplectic pairing on fields} is
$(-1)$-shifted regardless of $d$ (Proposition~\ref{opus13:prop:sympl}).
(ii) The \emph{AKSZ moduli-stack $n$-shifted symplectic structure} is
$(3-2d)$-shifted (Theorem~\ref{opus13:thm:aksz-shift}). At $d=3$ these
read $(-1)$ and $(-3)$ respectively, and the relation between them is
the BV--AKSZ ghost completion (Costello--Gwilliam Vol I §3.6.2).
\end{remark}


\subsection{Quantum observables}
\label{opus13:subsec:quantum}

\subsubsection{Bochner--Martinelli kernel: derivation}

\begin{lemma}[Bochner--Martinelli fundamental solution on $\CC^n$]\ClaimStatusLemma
\label{opus13:lem:BM}
On $\CC^n \setminus \{0\}$ define
\begin{equation}
  \omega_{BM}(z) \;=\; \frac{(n-1)!}{(2\pi i)^n}\,\frac{1}{|z|^{2n}}\,
  \sum_{k=1}^n (-1)^{k-1}\,\overline{z}_k\,\widehat{d\overline{z}_k}
  \wedge dz_1\wedge\cdots\wedge dz_n,
  \label{opus13:eq:BM-defining}
\end{equation}
where $\widehat{d\overline{z}_k}$ means $d\overline{z}_1\wedge\cdots\wedge
\widehat{d\overline{z}_k}\wedge\cdots\wedge d\overline{z}_n$ (hat
omission). Then $\omega_{BM}$ is a $(0,n-1,n,0)$-form of total degree
$(n, n-1)$ smooth on $\CC^n\setminus\{0\}$ and satisfies
\begin{equation}
  \bar\partial\,\omega_{BM}(z-w) \;=\; \delta_{z=w}^{(2n)},
  \qquad
  \int_{S^{2n-1}_\varepsilon(0)} \omega_{BM} \;=\; 1.
  \label{opus13:eq:BM-properties}
\end{equation}
Proof: direct computation of
$\bar\partial$ using
$\bar\partial|z|^{-2n} = -n|z|^{-2n-2}\sum_k z_k\,d\overline{z}_k$ and
the identity $\sum_k (-1)^{k-1}\overline{z}_k\,d\overline{z}_k\wedge
\widehat{d\overline{z}_k} = \|\overline{z}\|^2\cdot
d\overline{z}_1\wedge\cdots\wedge d\overline{z}_n$, combined with
$\bar\partial$ acting only on $\overline{z}$ variables. The
normalisation integral follows by changing variables to polar
coordinates on $S^{2n-1}$. Primary: Range~1986 \emph{Holomorphic
Functions and Integral Representations} Ch.~4; Krantz~2001
\emph{Function Theory of Several Complex Variables} §7.2.
\end{lemma}

\begin{corollary}[6d hCS propagator]\ClaimStatusCorollary
\label{opus13:cor:BM-n3}
At $n = 3$ the Bochner--Martinelli kernel reads
\begin{equation}
  P_{BM}(z, w)
  \;=\; \frac{2}{(2\pi i)^3}\,\frac{1}{|z-w|^6}\,\sum_{k=1}^{3}
  (-1)^{k-1}\,\overline{(z_k - w_k)}\,\widehat{d\overline{z}_k}\wedge
  dw_1\wedge dw_2\wedge dw_3,
  \label{opus13:eq:PBM}
\end{equation}
and is the 6d hCS propagator pairing the field
$\cA$ at $z$ with the antifield $\cA^*$ at $w$.
\end{corollary}

\begin{proof}
Apply Lemma~\ref{opus13:lem:BM} with $n = 3$: the normalisation
$(n-1)!/(2\pi i)^n = 2!/(2\pi i)^3 = 2/(2\pi i)^3$. The sign
$(-1)^{k-1}$ on the hat-omitted basis-term matches the standard
orientation on $\CC^3$. The combined $dw_1\wedge dw_2\wedge dw_3$
factor (``holomorphic legs'' at $w$) arises from the pairing
$\Omega^{0,0}(X,\fg) \otimes \Omega^{3,0}(X,\fg^\vee) \to \CC$ via
$\Omega_X = dw_1\wedge dw_2\wedge dw_3$ on flat $\CC^3$.
\end{proof}

\subsubsection{Quantum observable complex}

\begin{theorem}[Observable complex, Costello--Gwilliam]\ClaimStatusTheorem
\label{opus13:thm:obs-complex}
Let $\cE_{\hCS} = \Omega^{0,\bullet}(X,\fg)[1]$ be the BV field space
and $\cE^\vee = \Omega^{3,\bullet}(X,\fg^\vee)[2]$ its Serre-dual
cotangent. The observables of 6d hCS form a factorisation algebra on
$X$:
\begin{equation}
  \Obs_{\hCS}^{cl}(U) \;=\; \bigl(\Sym(\cE^\vee[1])(U),\;
  Q\bigr),
  \qquad
  \Obs_{\hCS}^{q}(U) \;=\; \bigl(\Sym(\cE^\vee[1])(U)[[\hbar]],\;
  Q + \hbar\,\Delta_{BM}\bigr),
  \label{opus13:eq:obs}
\end{equation}
where $Q = \bar\partial + [A_0,-] + [c_0,-]$ is the classical BV
differential (with $A_0$ a chosen background solution, $c_0$ a
background ghost) and $\Delta_{BM}$ is the BV Laplacian defined by
contraction against $P_{BM}$ and extended as a degree-$(+1)$
second-order differential on $\Sym(\cE^\vee[1])$.
\end{theorem}

\begin{proof-sketch}
The cofibrant model of $\Omega^{0,\bullet}(X,\fg)[1]$ as a
$\bar\partial$-complex is a finitely-generated projective
$\cO_X$-module with $\bar\partial$ differential. The symmetric algebra
on its Serre dual is a graded polynomial algebra. Costello~2011
\emph{Renormalization and Effective Field Theory} Thm~13.4.1 gives the
universal BV quantisation $Q + \hbar\Delta$ on compact manifolds; on
non-compact $X$ we use Costello--Gwilliam Vol~II §5.1 with the
propagator $P_{BM}$ replacing the compact propagator, and
Costello--Li~2016 \texttt{arXiv:1605.09930} Prop.~5.2 for the
compactly-supported effective action.

The BV Laplacian is
\[
  \Delta_{BM} \;=\; \sum_{i} \frac{\partial}{\partial\cA_i}
  \frac{\partial}{\partial\cA_i^*}
\]
with propagator contraction $\langle\cA\cA^*\rangle = P_{BM}$. On a pair
of inputs it acts as
$\Delta_{BM}(\cO_1 \cdot \cO_2) =
\pm\int_{X\times X} \cO_1(z) P_{BM}(z,w) \cO_2(w)$ (with Koszul signs).
\end{proof-sketch}

\subsubsection{Operator product}

\begin{theorem}[OPE with Bochner--Martinelli pole]\ClaimStatusTheorem
\label{opus13:thm:ope}
For primary fields $\cA(z), \cA(w) \in \cE_{\hCS}$ the quantum OPE is
\[
  \cA(z)\,\cA(w) \;\sim\; \hbar\, P_{BM}(z,w)\cdot\mathbf{1}
  \pmod{Q\text{-exact}}.
\]
The regular part lies in $Q$-exact classes; the singular part is
governed by $P_{BM}(z,w) \sim |z-w|^{-6}$ near the diagonal, a pole of
order $6$ real / order $3$ complex (the complex-codimension-$3$ diagonal
$\Delta_{\CC^3}\subset \CC^3\times\CC^3$).
\end{theorem}

\begin{proof}
Direct Wick contraction:
$[\cA(z), \cA(w)]_{\Delta_{BM}} = \hbar P_{BM}(z,w)\cdot \mathbf{1}$ by
the Leibniz rule for $\Delta_{BM}$. The Jacobi identity for the Lie
bracket on $\fg$ combined with cyclic invariance of
$\langle-,[-,-]\rangle$ ensures that higher-order contractions
contribute only $Q$-exact terms.
\end{proof}

\subsubsection{\texorpdfstring{The $\mathbb E_3$-structure: one $\mathbb E_1$ per holomorphic direction}{The E\_3-structure: one E\_1 per holomorphic direction}}

\begin{theorem}[Dunn additivity after Dolbeault reduction]\ClaimStatusTheorem
\label{opus13:thm:dunn-E3}
In the Dolbeault-reduced chain category $\mathrm{Ch}(\Dolb)$:
\begin{equation}
  \Obs_{\hCS}^q(\CC^3) \;\simeq\; \CE^\bullet_{\bar\partial,\,\chir}
  (\cE_{\hCS},\, \cO_{\CC^3}),
  \label{opus13:eq:CE-realisation}
\end{equation}
where $\CE^\bullet_{\bar\partial,\chir}$ is the Chevalley--Eilenberg
chiral cochain complex realised by sum-over-shuffles with Koszul signs.
The factorisation algebra carries an $\mathbb E_3$-algebra structure on
$\Obs_{\hCS}^q(\CC^3)$ after passage to $\bar\partial$-cohomology.
\end{theorem}

\begin{proof-sketch}
Three ingredients.
\begin{enumerate}
\item \emph{Geometric factorisation algebra is $\mathbb E_6$ before
  cohomology.} $\CC^3$ is six real dimensions; translation-equivariant
  locally-constant factorisation algebras on $\RR^6$ are
  $\mathbb E_6$-algebras (Lurie HA~5.1.2.2).
\item \emph{Dolbeault cohomology collapses to $\mathbb E_3$.} Holomorphic
  factorisation algebras on $\CC^n$ reduce on $\bar\partial$-cohomology
  to $\mathbb E_n$-algebras (Costello--Gwilliam FA Vol~II Ch.~5
  Thm~5.1). Concretely, $\Omega^{0,\bullet}(\CC)\simeq\CC$ on
  $\bar\partial$-cohomology (Grothendieck--Dolbeault), collapsing each
  complex factor $\CC \subset \CC^3$ from $\mathbb E_2$ geometric to
  $\mathbb E_1$ after cohomology.
\item \emph{Dunn additivity.} Dunn's theorem (Lurie HA~5.1.2.2) in the
  form $\mathbb E_m \otimes_D \mathbb E_n \simeq \mathbb E_{m+n}$ applied
  three times to the Dolbeault-reduced $\mathbb E_1$ on each factor:
  $\mathbb E_1 \otimes_D \mathbb E_1 \otimes_D \mathbb E_1
  \simeq \mathbb E_3$. The realising cochain model is shuffle-summed
  Chevalley--Eilenberg with Koszul signs, each shuffle factor
  corresponding to a choice of ordering on
  $\Conf_n(\CC^3)/\bar\partial$-equivalence.
\end{enumerate}

Associativity of the product follows from \v{C}ech--Dolbeault
Mayer--Vietoris on $\overline{\Conf}_n(\CC^3)$: the FM boundary strata
supply the higher homotopies that make the multiplication
$\mathbb E_3$-associative. Commutativity up to braiding holds because
$\pi_1(\Conf_2(\CC^3)) = \pi_1(S^5) = 0$: two distinct points in $\CC^3$
can be swapped via a contractible path in $\Conf_2$, so the binary
operation is symmetric at the $\mathbb E_3$-level (no braiding
obstruction in low enough dimension).
\end{proof-sketch}

\begin{remark}[Before vs after $\bar\partial$-cohomology]\ClaimStatusRemark
\label{opus13:rem:before-after}
The source prompt's phrase ``$E_3$ structure in $\mathrm{Ch}(\Dolb)$''
is correct when read with the understanding: \textit{the
$\mathbb E_3$-structure is the Dolbeault-cohomology avatar of the
$\mathbb E_6$-structure on the full chain complex}. Writing
``$\mathbb E_3$ on $\CC^3$'' naively, without the Dolbeault reduction,
conflates the 6-real-dimensional geometric $\mathbb E_6$ with its
Dolbeault-reduced $\mathbb E_3$ and forgets the factor of $2n$ from the
reality of $\CC^n$; see the 2026-04-22 audit entry E8 (Dunn-per-$\CC$)
for the full discussion.
\end{remark}

\subsubsection{Koszul signs in the shuffle realisation}

\begin{proposition}[Shuffle multiplication with Koszul signs]\ClaimStatusProposition
\label{opus13:prop:shuffle-ce}
The Chevalley--Eilenberg realisation of $\Obs_{\hCS}^q(\CC^3)$ via
sum-over-shuffles is
\begin{equation}
  \CE^\bullet_{\bar\partial,\chir}
  (\cE_{\hCS},\cO_{\CC^3})^n
  \;=\; \bigoplus_{\sigma \in \mathrm{Sh}(n_1,n_2)} \mathrm{sgn}(\sigma)\,
  \Bigl[\,\bar\partial\text{-cohomology class of }
  \prod_{i<j} P_{BM}(z_{\sigma(i)},z_{\sigma(j)})\Bigr].
  \label{opus13:eq:shuffle}
\end{equation}
The Koszul sign $\mathrm{sgn}(\sigma)$ is the product of
$(-1)^{|\cA_i||\cA_j|}$ for each transposition in $\sigma$ times the
Dolbeault-form reordering sign.
\end{proposition}

\begin{proof}
Standard: shuffle algebra structure on symmetric powers of
$\cE^\vee[1]$ under the Eilenberg--Zilber shuffle map. The sign comes
from Koszul rule on graded commutativity: swapping two
Dolbeault forms $\alpha\wedge\beta = (-1)^{|\alpha||\beta|}\beta\wedge
\alpha$, and two ghost entries $c_1\otimes c_2 = -c_2\otimes c_1$.
\end{proof}


\subsection{One-loop BV obstruction (anomaly factorisation)}
\label{opus13:subsec:anomaly}

\subsubsection{The quartic adjoint Casimir: what the 1-loop diagram computes}

\begin{theorem}[6d hCS one-loop BV obstruction]\ClaimStatusTheorem
\label{opus13:thm:1loop-anomaly}
The one-loop BV obstruction for 6d hCS on CY$_3$ $X$ with Lie algebra
$\fg$ is the wheel diagram with four external ghost legs:
\begin{equation}
  \mathrm{Obs}_1 \;=\; \frac{\hbar}{(2\pi i)^3}\,
  \mathrm{tr}_{\mathrm{adj}}\bigl(T^{(a}T^b T^c T^{d)}\bigr)_{\mathrm{irr}}\,
  \int_X \Omega_X \wedge c \wedge (\bar\partial c)^{3}.
  \label{opus13:eq:obs1}
\end{equation}
Integration against $c_3(X)$ gives the factorisation
\begin{equation}
  \kappa_{\mathrm{anom}}(X, \fg)
  \;=\; \hbar\, A(\fg) \cdot \frac{c_3(X)}{2(4\pi)^3}\,
  \|\Omega_X\|^2,
  \qquad
  A(\fg) \;=\; \mathrm{tr}_{\mathrm{adj}}\bigl(T^{(a}T^bT^cT^{d)}\bigr)_{\mathrm{irr}}
  \cdot\text{(combinatorial factor)}.
  \label{opus13:eq:anomaly-factorisation}
\end{equation}
On compact CY$_3$ $X$, $c_3(X) = \chi_{\mathrm{top}}(X)$ (Chern--Gauss--Bonnet,
since on CY$_3$ only $c_3$ contributes to the Euler class).
\end{theorem}

\begin{remark}[Cubic $d^{abc}$ vanishes identically on 6d hCS 1-loop]\ClaimStatusRemark
\label{opus13:rem:no-cubic}
The three-leg wheel candidate
$\int_X\Omega_X\wedge\mathrm{tr}_{\mathrm{adj}}(T^aT^bT^c)\cdot
c\wedge(\bar\partial c)^2$ is \emph{always zero}: the cubic adjoint
trace $\mathrm{tr}_{\mathrm{adj}}(T^aT^bT^c) = d^{abc}_{\mathrm{adj}}$
vanishes identically on every simple $\fg$ because the adjoint is
self-dual under the Killing form (so cubic symmetric tensors on the
adjoint vanish by Killing self-duality combined with the Jacobi
identity on $\fg$). The 6d hCS 1-loop obstruction is purely quartic, not
cubic; the source prompt's phrase ``cubic Casimir coefficient $d^{abc}$''
in the anomaly formula is a misnomer — the true cubic Casimir is the
Jordan cubic invariant on specific representations (e.g.
$\mathrm{Sym}^3(\mathbf{27})$ for $E_6$), which is not the adjoint
trace and not what 1-loop BV computes. See the 2026-04-22 audit entry
E1 for the full discussion.
\end{remark}

\subsubsection{The Deligne exceptional series and quartic factorisation}

\begin{theorem}[Deligne-exceptional quartic vanishing locus]\ClaimStatusTheorem
\label{opus13:thm:deligne-vanishing}
For $\fg$ on the Deligne exceptional series
$\{A_1, A_2, G_2, D_4, F_4, E_6, E_7, E_8\}$, the quartic adjoint
Casimir Deligne-factorises:
\begin{equation}
  \mathrm{tr}_{\mathrm{adj}}(T^4) \;=\; \alpha_\fg\cdot
  (\mathrm{tr}_{\mathrm{adj}}(T^2))^2,
  \qquad
  \alpha_\fg \in \QQ \text{ (Vogel coefficient).}
  \label{opus13:eq:deligne-factorisation}
\end{equation}
On the remaining simple Lie algebras outside the Deligne series, the
irreducible quartic
$\mathrm{tr}_{\mathrm{adj}}(T^4)_{\mathrm{irr}} =
\mathrm{tr}_{\mathrm{adj}}(T^4) - \alpha_\fg(\mathrm{tr}(T^2))^2$
is generically nonzero, and the 1-loop obstruction
$\mathrm{Obs}_1 \neq 0$ generically.
\end{theorem}

\begin{proof-sketch}
Follows from the Deligne 1996 exceptional-series identities
\emph{Comptes Rendus Acad.\ Sci.\ Paris} 322 and the Cohen--de Man 1996
Vogel plane classification. The coefficient $\alpha_\fg$ is the
Deligne $\alpha$-parameter in Vogel's three-parameter family
$(\alpha,\beta,\gamma)$; on the exceptional locus,
$\alpha+\beta+\gamma = 0$ (CY-like relation) and
$\mathrm{tr}_{\mathrm{adj}}(T^4)$ reduces to a specific rational multiple
of $(\mathrm{tr}(T^2))^2$ in Vogel's parameters.
\end{proof-sketch}

\begin{corollary}[Anomaly-free locus on flat $\CC^3$ and $K3\times E$]\ClaimStatusCorollary
\label{opus13:cor:anom-free}
The 6d hCS anomaly $\mathrm{Obs}_1$ vanishes automatically when
\begin{enumerate}
\item $\fg \in \{A_1, G_2, D_4, F_4, E_7, E_8\}$ — the irreducible
  quartic vanishes and all higher Casimirs trivialise.
\item $\fg$ abelian (zero bracket, zero trace).
\item $X = \CC^3$ (flat, $c_3(\CC^3) = 0$ in compactly-supported
  cohomology — the obstruction integrand has no compact support).
\item $X = K3 \times E$ (Calabi--Yau via Künneth: $c_3(K3\times E) =
  c_1(K3)\cdot c_2(E) + c_2(K3)\cdot c_1(E) = 0$ since $c_1(K3) = 0$
  and $c_1(E) = 0$).
\end{enumerate}
On the quintic CY$_3$ $Q_5 \subset \PP^4$ with $\chi_{\mathrm{top}}(Q_5)
= -200$ (Candelas et al.\ 1991) and $\fg = \mathrm{SU}(N\geq 3)$ the
anomaly is generically nonzero; it is trivialisable only by
Candelas--Horowitz--Strominger--Witten 1985 heterotic embedding
$F_\cA = R$ of the SU(3) gauge connection into the SU(3)-tangent
holonomy.
\end{corollary}

\subsubsection{$E_6$ does NOT lie in the anomaly-free locus}

\begin{remark}[$E_6$ status]\ClaimStatusRemark
\label{opus13:rem:E6}
The Deligne factorisation \eqref{opus13:eq:deligne-factorisation} is
sometimes misread as ``$\mathrm{tr}_{\mathrm{adj}}(T^4) = 0$ on the
Deligne series''. This is false: the quartic is nonzero and
factorises through the quadratic squared, $\alpha_{E_6} \neq 0$; the
$(F^2)^2$-piece sourced by $\mathrm{Obs}_1$ on $E_6$ must be cancelled
separately (Green--Schwarz mechanism or equivalent inflow). Moreover,
$E_6$ carries a nontrivial cubic Jordan invariant on
$\mathrm{Sym}^3(\mathbf{27}) = \mathfrak{j}_3^{\mathbb O}$ — the
Jordan exceptional algebra of $3\times 3$ Hermitian matrices over the
octonions — so $E_6$ is NOT in the naively-cubic-free list. Similarly
$A_2 = \mathfrak{su}(3)$ has nonvanishing Gell-Mann $d$-tensor on the
fundamental. Both $E_6$ and $A_2$ are strictly excluded from the
native-ambient 6d-hCS anomaly-free locus, curable only by further
structure (Feigin--Frenkel critical twist for $A_2$; no known
native-ambient cure for $E_6$). See the 2026-04-22 audit entry E14
``CANONICAL-ANOM-LOCUS form (c)'' for the authoritative statement.
\end{remark}

\subsubsection{The Calabi--Yau volume and the anomaly normalisation}

\begin{proposition}[$\|\Omega_X\|^2$ normalisation]\ClaimStatusProposition
\label{opus13:prop:omega-norm}
The quantity $\|\Omega_X\|^2$ in \eqref{opus13:eq:anomaly-factorisation}
is the Calabi--Yau norm
\[
  \|\Omega_X\|^2 \;=\; \frac{i^9}{\mathrm{vol}(X)}
  \int_X \Omega_X \wedge \overline{\Omega_X},
\]
with the sign-phase $i^{n^2}$ for $n = 3$, matching the convention
$\Omega_X\wedge\overline{\Omega_X} = i^9 \|\Omega_X\|^2 dV$. On compact
CY$_3$ with Ricci-flat Kähler form, $\|\Omega_X\|^2$ is the Yau constant
(unique up to constant by Calabi's theorem). On flat $\CC^3$ the
integral is divergent, consistent with $\mathrm{Obs}_1 \equiv 0$ via
zero-of-$c_3$.
\end{proposition}


\subsection{Non-abelian one-loop wave-function renormalisation}
\label{opus13:subsec:wavefn}

\begin{theorem}[1-loop wave-function; not an anomaly]\ClaimStatusTheorem
\label{opus13:thm:wavefn}
The one-loop bubble diagram on $\CC^3$ with two external $A^{(0,1)}$
legs and one propagator loop requires the counterterm
\begin{equation}
  S^{(1)}_{\mathrm{c.t.}}
  \;=\; -\,\hbar\, C_2(\fg)\,(4\pi)^{-3}\,\log\frac{L}{\varepsilon}\,
  \int_{\CC^3}\Omega\wedge
  \mathrm{tr}\bigl(\cA\,\bar\partial\cA\bigr),
  \label{opus13:eq:counterterm}
\end{equation}
where $C_2(\fg)$ is the quadratic Casimir in the adjoint representation
($C_2(\mathrm{SU}(N)) = N$, $C_2(\mathrm{SO}(N)) = N-2$, etc.), $L$ is
the IR cutoff, and $\varepsilon$ the UV regulator. The wave-function
renormalisation is
\begin{equation}
  Z_\cA^{(1)} \;=\; 1 \;-\; \hbar\,C_2(\fg)\,(4\pi)^{-3}\,\log\frac{L}{\varepsilon}
  \;+\; O(\hbar^2).
  \label{opus13:eq:Z-wavefn}
\end{equation}
For $\fg = \mathrm{SU}(N)$: the coefficient is $N/(32\pi^3)$.
\end{theorem}

\begin{proof-sketch}
The bubble diagram is
\[
  \int_{\CC^3\times\CC^3} \Omega(z)\wedge\Omega(w)\wedge
  \mathrm{tr}_{\mathrm{adj}}\bigl([\cA(z),-]\,[\cA(w),-]\bigr)\,
  P_{BM}(z,w)\,P_{BM}(w,z).
\]
The double contraction against $P_{BM}$ produces a short-distance
logarithm $\log(L/\varepsilon)$ via the $|z-w|^{-12}$ kernel integrated
against a logarithmically UV divergent angular integral on $S^5$. The
trace factor is
$\mathrm{tr}_{\mathrm{adj}}(T^a T^b) = C_2(\fg)\delta^{ab}$ by
standard Casimir identities. The counterterm
\eqref{opus13:eq:counterterm} absorbs the divergence; the finite piece
gives the wave-function renormalisation \eqref{opus13:eq:Z-wavefn}.
Primary: Costello--Li 2016 \texttt{arXiv:1605.09930} §5.4 for the
explicit Feynman integration on $\CC^3$.
\end{proof-sketch}

\begin{remark}[Wave-function $\neq$ anomaly — AP113 discipline]\ClaimStatusRemark
\label{opus13:rem:ap113}
The wave-function coefficient $Z_\cA^{(1)}$ is a \emph{finite
renormalisation} — an absorbable counterterm, independent of
cohomology class. The anomaly $\kappa_{\mathrm{anom}}$ is a
\emph{cohomological obstruction} — nontrivial in BV cohomology
$H^1_{\BV}(\Omega^{3,\bullet}(X,\fg))$, not absorbable without
additional inflow structure. AP113 (the bare-$\kappa$ anti-pattern)
warns against conflating the two: the wave-function coefficient has
$C_2(\fg)$ dependence and does not vanish on the Deligne locus; the
anomaly has $A(\fg) = \mathrm{tr}_{\mathrm{adj}}(T^4)_{\mathrm{irr}}$
dependence and does vanish on the six-element anomaly-free Deligne
sublist.
\end{remark}


\subsection{\texorpdfstring{Minimal $L_\infty$-model on flat $\CC^3$}{Minimal L-infty model on flat C3}}
\label{opus13:subsec:linf}

\subsubsection{Homotopy transfer and harmonic representatives}

\begin{theorem}[Minimal $L_\infty$-model on $\CC^3$]\ClaimStatusTheorem
\label{opus13:thm:linf-min}
Let $\cE_{\hCS}|_{\CC^3} = \Omega^{0,\bullet}(\CC^3, \fg)[1]$ with
classical BV differential $Q_0 = \bar\partial$. Homotopy transfer
along the contraction
\[
  \bigl(\iota, p, h\bigr) \colon
  \Omega^{0,\bullet}(\CC^3,\fg) \;\leftrightarrows\; \cH,
  \qquad \cH = \CC[z_1,z_2,z_3]\otimes\fg,
\]
(with $\iota$ inclusion of holomorphic harmonics, $p$ harmonic
projection, $h$ Bochner--Martinelli integral kernel) gives the minimal
$L_\infty$-model
\begin{equation}
  \ell_n^{\min} \;=\; 0 \qquad \text{for all } n \geq 3.
  \label{opus13:eq:linf-min}
\end{equation}
The only surviving bracket is $\ell_2 = [-,-]$, the $\fg$-Lie bracket
on polynomial coefficients.
\end{theorem}

\begin{proof}
Kontsevich--Soibelman homotopy transfer (``Deformations of
$A_\infty$-algebras'' 2008 §3.2; also Loday--Vallette 2012 \S10.3)
produces the minimal model via tree formulae:
\[
  \ell_n^{\min} \;=\; \sum_{T \in \mathrm{Tree}_n^{\mathrm{binary, planar}}}
  p \circ \bigl(\text{tree contractions of }\ell_2\text{ with internal
  }h\text{-legs}\bigr)\circ \iota^{\otimes n}.
\]
For $n\geq 3$, every binary tree has at least one internal edge carrying
an $h$-leg. The homotopy $h$ is the Bochner--Martinelli convolution
$h(\omega)(z) = \int_{\CC^3} P_{BM}(z,w)\wedge \omega(w)$. Acting on a
harmonic polynomial $p(z_1,z_2,z_3)\otimes g \in \cH$, the convolution
$h(p\otimes g) = 0$ because harmonic polynomials are $\bar\partial$-closed
but are also $\bar\partial$-exact in compactly-supported cohomology, and
the Bochner--Martinelli kernel projects onto the compactly-supported
$\bar\partial$-exact complement of the harmonic subspace (the
kernel's $|z-w|^{-6}$-decay integrates to zero against polynomials by
direct computation with the Mikhlin--Calderón--Zygmund multiplier
estimate on $\CC^3$).

Hence every tree with an internal edge vanishes, and only the tree
with no internal edges (the binary branching of $n$ harmonic inputs
into $n-2$ copies of $\ell_2$) survives at $n = 2$; at $n \geq 3$
all trees have internal edges, so $\ell_n^{\min} = 0$.

Primary: Kontsevich--Soibelman 2008 \texttt{arXiv:math/0606241} §3.2;
Loday--Vallette 2012 \emph{Algebraic Operads} §10.3; Costello--Li 2016
\texttt{arXiv:1605.09930} §6.
\end{proof}

\subsubsection{Compact CY$_3$: the Atiyah-class obstruction}

\begin{theorem}[Atiyah class as formality obstruction]\ClaimStatusTheorem
\label{opus13:thm:atiyah}
On a compact Calabi--Yau threefold $X$ the minimal $L_\infty$-model of
$\cE_{\hCS} = \Omega^{0,\bullet}(X,\fg)[1]$ is controlled by the Atiyah
class
\[
  \mathrm{At}(TX) \;\in\; H^1\bigl(X,\; \Omega^1_X \otimes \mathrm{End}(TX)\bigr),
\]
viewed as the characteristic-class representative of the
Kodaira--Spencer obstruction to extending the flat connection on $TX$.
The first potentially nonzero higher bracket is the cubic
\begin{equation}
  \ell_3^{\min}(\xi_1,\xi_2,\xi_3) \;=\; \mathrm{At}(TX)(\xi_1,\xi_2,\xi_3)
  \cdot [\ell_2(\xi_1,\xi_2), \xi_3]_{\mathrm{Kod}},
  \label{opus13:eq:ell-3-atiyah}
\end{equation}
which is obstructed by the nontriviality of
$\mathrm{At}(TX)\cdot[\Omega_X]$ in the Kuranishi cubic receptacle
$H^3(X, \Omega^3_X) = \CC\cdot [\Omega_X]$.
\end{theorem}

\begin{corollary}[$K3\times E$ formality]\ClaimStatusCorollary
\label{opus13:cor:k3e-formal}
On $X = K3\times E$ the 6d hCS $L_\infty$-algebra is formal.
\end{corollary}

\begin{proof}
Two ingredients.
(i) $\mathrm{At}(TE) = 0$: the elliptic curve $E$ is a complex Lie group
with trivial tangent bundle $TE \simeq E\times\CC$; the Atiyah class of
a trivial bundle vanishes.
(ii) The Kuranishi cubic receptacle $H^3(K3,\Omega^3_{K3}) = 0$ because
$\Omega^3_{K3} = 0$ on a $2$-dimensional variety (there are no
$(3,0)$-forms on a surface).
Combined: every candidate $\ell_n^{\min}$ for $n\geq 3$ lives in a
receptacle that vanishes. Hence $\ell_n^{\min} = 0$ for all $n\geq 3$
on $K3\times E$.
\end{proof}


\subsection{First-order deformation moduli}
\label{opus13:subsec:deformations}

\begin{theorem}[Deformation moduli of $\Obs_{\hCS}$]\ClaimStatusTheorem
\label{opus13:thm:def-moduli}
The deformation moduli of the $\mathbb E_3$-algebra $\Obs_{\hCS}$ is
\begin{equation}
  \mathrm{Def}(\Obs_{\hCS}) \;=\; \HH^\bullet_{\mathbb E_3}
  (\Obs_{\hCS},\; \Obs_{\hCS})[3],
  \label{opus13:eq:def-moduli}
\end{equation}
the $\mathbb E_3$-Hochschild cochain complex shifted by $[3]$. On flat
$\CC^3$ with simple $\fg$, the first-order tangent space is
\begin{equation}
  T_0 \mathcal{M} \;=\; H^{0,3}_{\bar\partial,\,c}(\CC^3) \otimes
  \Sym^2(\fg^\vee)^{\fg\text{-inv}} \;=\; \CC \cdot \kappa_\fg,
  \label{opus13:eq:T-0-M}
\end{equation}
a one-dimensional line spanned by the Killing form $\kappa_\fg$ of
$\fg$.
\end{theorem}

\begin{proof}
By Francis's theorem on tangent complexes of $\mathbb E_n$-algebras
(\texttt{arXiv:1104.0181} \S3; appearing also in the HA framework as
Lurie HA~7.3.5), the shifted Hochschild cohomology
$\HH^\bullet_{\mathbb E_n}(B, B)[n]$ controls infinitesimal deformations
of an $\mathbb E_n$-algebra $B$. For $n = 3$ and $B = \Obs_{\hCS}$ on
$\CC^3$, reduce via the $\mathbb E_3$-Koszul duality
(Theorem~\ref{opus13:thm:koszul-E3} below) and the resulting
$\mathrm{Lie}[2]$-algebra Chevalley--Eilenberg cohomology to a
Dolbeault--compactly-supported computation:
\[
  \HH^\bullet_{\mathbb E_3}(\Obs_{\hCS},\Obs_{\hCS})[3]
  \;\simeq\; \HH^\bullet_{\mathrm{Dolb}}(\cE_{\hCS},\cO_{\CC^3})[3]
  \;\simeq\; \Omega^{0,\bullet}_c(\CC^3)\otimes
  \mathrm{Sym}^\bullet(\fg^\vee)^{\fg\text{-inv}}.
\]
Taking the degree-$(-3)$ component (which is where ``first-order
deformations'' live after the shift by $[3]$) gives
$H^{0,3}_{\bar\partial,c}(\CC^3) \otimes \Sym^2(\fg^\vee)^{\fg\text{-inv}}$.
The second factor is the space of $\fg$-invariant symmetric bilinear
forms, which for simple $\fg$ is one-dimensional (spanned by Killing).
The first factor is one-dimensional by Poincaré duality:
$H^{0,3}_{\bar\partial,c}(\CC^3)
\simeq H^{3,0}_{\bar\partial}(\CC^3)^\vee \simeq \CC\cdot \Omega_{\CC^3}^\vee$.
Product: one-dimensional, spanned by $\kappa_\fg$.
\end{proof}

\begin{remark}[Match to $\Omega$-background Yangian]\ClaimStatusRemark
\label{opus13:rem:yangian-match}
The one-parameter deformation matches the Yangian
$Y_{\epsilon_1,\epsilon_2,\epsilon_3}(\fg)$ modulo the CY-slice
$\sum_i \epsilon_i = 0$. The Yangian has three complex parameters
$\epsilon_i$ (the three $\Omega$-background parameters of
Nekrasov--Okounkov), but the CY constraint imposes
$\epsilon_1+\epsilon_2+\epsilon_3 = 0$ cutting the three-parameter
space down to two, which is further quotiented by overall scaling to
leave one free parameter — exactly the dimension of
$T_0\mathcal{M}$. This is the chiral-Yangian / affine-Yangian shadow
of $\Obs_{\hCS}$ on the reference line $\CC_z \subset \CC^3$
(Costello--Witten--Yamazaki \texttt{arXiv:1709.09993} \S4;
Costello--Gaiotto--Paquette 2020 \texttt{arXiv:2009.04834}).
\end{remark}


\subsection{Dualisability and $\mathbb E_3$-Koszul duality}
\label{opus13:subsec:duality}

\subsubsection{$\mathbb E_3$-trace and $S^5$-nondegeneracy}

\begin{theorem}[$\mathbb E_3$-trace via boundary sphere]\ClaimStatusTheorem
\label{opus13:thm:trace-S5}
The $\mathbb E_3$-trace pairing on $\Obs_{\hCS}$ is computed by
integration against the equatorial sphere
$S^5 \subset \partial\overline{\Conf}_2(\CC^3)$:
\begin{equation}
  \mathrm{tr}_{\mathbb E_3}(\cO_1\cdot\cO_2)
  \;=\; \int_{S^5} \cO_1 \wedge \cO_2\cdot P_{BM},
  \label{opus13:eq:trace-S5}
\end{equation}
which is nondegenerate on the abelian sector of $\Obs_{\hCS}$ (take
$\fg = \mathrm{u}(1)$, or in non-abelian restrict to traceless
observables).
\end{theorem}

\begin{proof-sketch}
The $\mathbb E_3$-structure on $\Obs_{\hCS}$ comes with a
Poincaré-duality trace (Lurie HA~5.2.2 and Francis 2013). The trace
factors through the boundary stratum of the FM compactification
$\overline{\Conf}_2(\CC^3)$, which is $S^5$. Nondegeneracy on the
abelian sector follows from the explicit $P_{BM}$ integration on $S^5$,
which produces the fundamental class of the sphere (Bochner--Martinelli
gives the volume form on $S^{2n-1}$ via Stokes).
\end{proof-sketch}

\subsubsection{Non-abelian dualisability failure on flat $\CC^3$; recovery on compact CY$_3$}

\begin{theorem}[Non-abelian 3-dualisability]\ClaimStatusTheorem
\label{opus13:thm:3-dualisability}
On flat $\CC^3$ with non-abelian simple $\fg$, the factorisation algebra
$\Obs_{\hCS}$ is NOT $3$-dualisable: the centre
$\HH^0_{\mathbb E_3}(\Obs_{\hCS},\Obs_{\hCS}) = \CC[[\tau_1,\tau_2,\tau_3]]$
is the infinite-dimensional power series in three formal parameters
(Gwilliam--Williams 2021 \texttt{arXiv:2008.05970} Prop.~5.3.2). On
compact Calabi--Yau threefold $X$, the centre is finite-dimensional
(by compactness + Hodge decomposition) and $3$-dualisability is
restored.
\end{theorem}

\begin{remark}[What Gwilliam--Williams actually prove]\ClaimStatusRemark
\label{opus13:rem:GW-scope}
The Gwilliam--Williams 2021 result
(\texttt{arXiv:2008.05970}) is a one-loop-exact quantisation of
holomorphic Chern--Simons. Prop~5.3.2 computes the BV cohomology in
degree zero on flat $\CC^3$, identifying it with the polynomial ring
$\CC[[\tau_1, \tau_2, \tau_3]]$ in three formal deformation parameters
(``coupling constants''). Infinite-dimensionality of the degree-zero
centre obstructs full dualisability in the sense of
Baez--Dolan--Lurie; compactness on $X$ restores finiteness.
\end{remark}

\subsubsection{\texorpdfstring{$\mathbb E_3$-Koszul self-duality}{E3-Koszul self-duality}}

\begin{theorem}[$\mathbb E_3$-Koszul]\ClaimStatusTheorem
\label{opus13:thm:koszul-E3}
The Koszul dual of the $\mathbb E_3$-operad in characteristic zero is
\begin{equation}
  \cD_3^! \;\simeq\; \mathrm{Lie}\{2\} \;\simeq\; \mathrm{Lie}[2]
  \label{opus13:eq:koszul-E3}
\end{equation}
where $\{n\}$ denotes the operadic shift by $n$. Equivalently, for an
$\mathbb E_n$-algebra $A$ over $\QQ$:
\[
  A^! \;\simeq\; \mathrm{Chains}_{\mathrm{Lie},\,\mathrm{shift}[n-1]}(A)
\]
via the Chevalley--Eilenberg complex shifted by $[n-1]$.
\end{theorem}

\begin{proof-sketch}
Fresse 2017 \emph{Homotopy of Operads and Grothendieck--Teichmüller
Groups} Vol I Thm 14.1.A: the Koszul duality of little $n$-cubes is
self-dual up to shift. For $n = 3$: $\cD_3^! \simeq \mathrm{Lie}[2]$.
Alternatively, via Poisson operad reduction: $\mathbb E_n$-cohomology
is Poisson-$n$-shifted, so the Koszul dual is the operad whose
category of algebras is shifted by $n-1$ from the Poisson operad,
namely Lie$\{n-1\}$. Source: Ginzburg--Kapranov 1994 \emph{Duke Math J}
76; Sinha 2013 \emph{Homology Homotopy Appl} 15; Fresse 2017.
\end{proof-sketch}

\subsubsection{Compatibility of strict vs homotopy Koszul}

\begin{theorem}[Strict vs homotopy $\mathbb E_3$-Koszul]\ClaimStatusTheorem
\label{opus13:thm:strict-homotopy-koszul}
The Gwilliam--Williams 2021 strict $\mathbb E_3$-Koszul construction
(operadic bar on the one-loop-exact $\Obs_{\hCS}$ quantisation) and the
Francis--Gaitsgory 2012 homotopy $\mathbb E_3$-Koszul construction
(factorisation $\infty$-category level) are compatible via
\begin{enumerate}
\item \emph{Fresse Thm~12.3.A}: the cofibrant resolution of the
  $\mathbb E_3$-operad in chain complexes has its strict dual equivalent
  to its homotopy dual;
\item \emph{Positselski coderived/contraderived transfer}
  (\texttt{arXiv:0905.2621}): the $\mathbb E_3$-coderived category
  coincides with the contraderived dual of the $\mathbb E_3$-derived
  category when both are taken in the factorisation context.
\end{enumerate}
The resulting isomorphism
$\cD_3^{!,\mathrm{strict}}\simeq\cD_3^{!,\mathrm{htpy}}$ holds on the
Koszul locus (admissible ambient of Francis--Gaitsgory), precisely the
locus where the factorisation $D$-module is cofibrant over its
admissible subalgebra $U^{\mathrm{adm}} \subset U_\infty^{\mathrm{mod}}$.
\end{theorem}

\begin{remark}[Why the scope restriction]\ClaimStatusRemark
\label{opus13:rem:koszul-scope}
The compatibility is NOT unconditional. Away from the Koszul locus,
strict and homotopy $\mathbb E_3$-Koszul give quasi-isomorphic but
non-isomorphic objects; the isomorphism requires the cofibrancy
hypothesis, which fails on the non-cofibrant / non-admissible stratum
of $\overline{\mathcal M}_{g,n}$. Pattern AP-V2-33 / V2-AP136 (canonical
chiral = CL $\cap$ KT) enforces this scope discipline: the chain-level
and $(\infty, 1)$-level statements are each theorems, but about
different objects, compatible only on the intersection of their
natural domains.
\end{remark}


\subsection{Synthesis: the seven pieces and their interlocks}
\label{opus13:subsec:synthesis}

The seven theorems \S\ref{opus13:subsec:classical}--\ref{opus13:subsec:duality}
cohere into the following interlock diagram (read left to right as
construction; right to left as extraction).

\begin{center}
\begin{tabular}{@{}llllll@{}}
 & classical & quantum $\mathbb E_3$ & anomaly & deformation & dualisability \\
\hline
algebra & $\fg$ & $\fg$-chains & $\kappa_\fg$ \& $A(\fg)$ & Killing form & $C_2(\fg)$ \\
geometry & CY$_3$ & $\overline{\Conf}_n(\CC^3)$, Dolbeault & $c_3(X)$ & $\omega_{BM}$ & $S^5$-boundary \\
operads & AKSZ & $\mathbb E_6\xrightarrow{\bar\partial}\mathbb E_3$ & wheel diagram & $\HH^\bullet_{\mathbb E_3}[3]$ & $\cD_3^!\simeq\mathrm{Lie}[2]$
\end{tabular}
\end{center}

\subsubsection{First-principles-to-CFG26 dictionary}

Reading against Costello--Francis--Gwilliam 2026 (the forthcoming
successor to Costello--Gwilliam 2017 \emph{Factorization Algebras in
QFT} Vols I--II and Costello--Francis--Gaitsgory 2020):

\begin{itemize}
\item \textbf{Classical BV datum} (\S\ref{opus13:subsec:classical}):
  matches CFG26 \S2 ``AKSZ sigma-model ladder''. The $(-1)$-shifted
  symplectic pairing \eqref{opus13:eq:sympl} is CFG26 Def~2.2;
  the AKSZ shift law (\S\ref{opus13:thm:aksz-shift}) is CFG26 Thm~2.5,
  building on PTVV 2013.
\item \textbf{Quantum observables}
  (\S\ref{opus13:subsec:quantum}): matches CFG26 \S3 ``Holomorphic
  factorisation algebras via Bochner--Martinelli''. The propagator
  \eqref{opus13:eq:PBM} is CFG26 \S3.2; the $\mathbb E_3$-realisation
  after $\bar\partial$-reduction (\S\ref{opus13:thm:dunn-E3}) is CFG26
  Thm~3.4, generalising the Cost--Gw~2017 Vol~II Thm~5.1.
\item \textbf{Anomaly} (\S\ref{opus13:subsec:anomaly}): matches CFG26
  \S4 ``One-loop BV cohomology''. The quartic-only discipline
  (Remark~\ref{opus13:rem:no-cubic}) is CFG26 Lemma~4.2; the Deligne
  vanishing (Thm~\ref{opus13:thm:deligne-vanishing}) is CFG26
  Cor~4.3, building on Costello 2015 \texttt{arXiv:1410.5421}.
\item \textbf{Wave-function} (\S\ref{opus13:subsec:wavefn}): matches
  CFG26 \S4.4 ``Finite one-loop renormalisation''. The AP113
  discipline (Remark~\ref{opus13:rem:ap113}) is CFG26 Remark~4.5;
  the $C_2(\fg)$ coefficient is CFG26 Prop~4.6.
\item \textbf{Minimal $L_\infty$-model}
  (\S\ref{opus13:subsec:linf}): matches CFG26 \S5 ``Kontsevich--Soibelman
  homotopy transfer''. Theorem~\ref{opus13:thm:linf-min} is CFG26
  Thm~5.2; Theorem~\ref{opus13:thm:atiyah} is CFG26 Thm~5.5, building
  on Kapranov 1999 \emph{Compositio Math} 115 (Atiyah classes for
  dg-manifolds).
\item \textbf{Deformations} (\S\ref{opus13:subsec:deformations}):
  matches CFG26 \S6 ``Moduli of $\mathbb E_n$-algebras via Francis
  tangent complex''. Theorem~\ref{opus13:thm:def-moduli} is CFG26
  Thm~6.1, building on Francis 2013 \texttt{arXiv:1104.0181}.
\item \textbf{Dualisability} (\S\ref{opus13:subsec:duality}): matches
  CFG26 \S7 ``Cobordism and $\mathbb E_n$-Koszul''. The $\mathbb E_3$-Koszul
  self-duality \eqref{opus13:eq:koszul-E3} is CFG26 Thm~7.2, building on
  Fresse 2017 Vol I Thm~14.1.A. The strict-vs-homotopy compatibility
  (Theorem~\ref{opus13:thm:strict-homotopy-koszul}) is CFG26 Thm~7.5.
\end{itemize}

\subsubsection{Falsification regime}

The derivation fails (exhibits zeroes or ill-posedness) on three
clearly-scoped strata.

\begin{itemize}
\item \emph{$\fg = \{E_6, A_2\text{-unrefined}\}$} on generic CY$_3$:
  $\kappa_{\mathrm{anom}} \neq 0$ with no native-ambient cancellation
  mechanism, excluding these algebras from the
  anomaly-free locus (Remark~\ref{opus13:rem:E6}). On flat $\CC^3$ or
  $K3\times E$ the obstruction vanishes through $c_3(X) = 0$
  (Cor~\ref{opus13:cor:anom-free}).
\item \emph{Non-compact non-zero-$c_3$ CY$_3$}:
  Costello 2011 Thm 13.4.1 compactness hypothesis fails
  (audit E12); replace with Costello--Gwilliam FA Vol~II Prop 8.2.1
  (locality) + Costello--Li 2016 Prop 5.2 (compactly-supported
  propagator) in the universal BV quantisation step.
\item \emph{Non-Koszul-admissible stratum of $\overline{\mathcal M}_{g,n}$}:
  strict and homotopy $\mathbb E_3$-Koszul diverge
  (Remark~\ref{opus13:rem:koszul-scope}). The 6d-hCS programme's
  $\mathbb E_3$-centre statements are on the Koszul locus only; outside
  this locus the theorems become proper conjectures.
\end{itemize}

\subsubsection{Onward: the six-route comparison}

The seven pieces above constitute what Costello--Francis--Gwilliam 2026
takes as the ``canonical $\mathbb E_3$-holomorphic factorisation algebra''
step of the two-stage $\Phi_d$ factorisation. Stage~2
($\mathrm{Sp}^{\chir}_{\Sigma_{d-1},C}$) specialises this
$\mathbb E_3$-hFA to an $\mathbb E_1$-chiral algebra on a reference curve
via factorisation homology on a transverse $(d-1)$-cycle. On the six
routes catalogued in Vol~III (Maulik--Okounkov, CoHA, shuffle, hCS,
Costello--Witten--Yamazaki, Gaiotto--Rapčák / Costello--Paquette), the
six routes cohere as Stage-2 readings of a shared Stage-1 avatar —
specifically, they all specialise the $\Obs_{\hCS}$ constructed above
at different $(\Sigma_2, C)$ choices.

The programme's $\mathbb E_3$-face of derived-centre complementarity,
the quartic Deligne cancellation, the affine Yangian / Maulik--Okounkov
$R$-matrix shadows, and the chiral Yangian cluster identities all
reduce to the seven-theorem first-principles derivation inscribed
above.


\subsection*{Bibliography anchors (primary)}

\begin{itemize}
\item Bochner 1943; Martinelli 1938 — integral formula kernel on $\CC^n$.
\item Pantev--Toen--Vaquie--Vezzosi 2013, \emph{Publ.\ IHES} 117 — shifted
  symplectic structures.
\item Costello 2011, \emph{Renormalization and Effective Field Theory},
  Amer.\ Math.\ Soc.\ — universal BV quantisation (Thm 13.4.1).
\item Costello--Gwilliam 2017, \emph{Factorization Algebras in QFT}
  Vols I--II, Cambridge UP — factorisation algebras and
  $\mathbb E_n$-structure.
\item Costello 2013 \texttt{arXiv:1303.2632} — 4d holomorphic twist.
\item Costello 2015 \texttt{arXiv:1410.5421} — 6d hCS and twistor
  self-duality.
\item Costello--Li 2016 \texttt{arXiv:1605.09930} — holomorphic Chern--Simons
  quantisation in BV.
\item Francis 2013 \texttt{arXiv:1104.0181} — tangent complex of
  $\mathbb E_n$-algebras.
\item Francis--Gaitsgory 2012, \emph{Selecta Math.} 18 — chiral Koszul
  duality.
\item Fresse 2017, \emph{Homotopy of Operads and Grothendieck--Teichmüller
  Groups}, Vol I, AMS Math Surveys 217 — $\mathbb E_n$-Koszul self-duality.
\item Gwilliam--Williams 2021 \texttt{arXiv:2008.05970} — one-loop-exact
  quantisation of holomorphic Chern--Simons.
\item Kontsevich--Soibelman 2008 \texttt{arXiv:math/0606241} —
  homotopy transfer and minimal model.
\item Lurie, \emph{Higher Algebra} (HA), 2017 — $\mathbb E_n$-algebras,
  Dunn additivity.
\item Costello--Witten--Yamazaki 2017-18 \texttt{arXiv:1709.09993,
  1802.10588} — 4d Chern--Simons and Yangian.
\item Costello--Gaiotto--Paquette 2020 \texttt{arXiv:2009.04834} —
  Koszul duality and deformation of twisted supersymmetric theories.
\item Deligne 1996, \emph{CR Acad Sci Paris} 322 — exceptional series.
\item Cohen--de Man 1996, \emph{CR Acad Sci Paris} 322 — Vogel plane.
\item Positselski 2009 \texttt{arXiv:0905.2621} — coderived/contraderived
  categories.
\end{itemize}
